\scanpage{180}
\[
f^\circ(\alpha_1,u_2(x_2),\ldots,u_n(x_n)).
\]
\margin{\uline{NB} Si les \deleted{\unsure{ext.}} $E_i$ sont munis de scindages et leurs structures de $\lambda$-triv. définies par des $\unsure{b_i}:F_i\unsure{\to|-}E_i^\circ$, alors la $\lambda$-trivialisation de $\mathop{*}E_i$, scindé de la façon \unsure{can.}, est donnée par $\bigotimes b_i:\bigotimes F_i\unsure{\to|-}\bigotimes E_i^\circ$.}
Ainsi, on a une opération interne $*$ dans la catégorie des extensions $\lambda$-scindées \inserted{(à quotients $F_i$ plats, au besoin)}, et en particulier dans celle des torseurs $\lambda$-scindés (car si les $F_i$ sont $\simeq k$, $\bigotimes F_i\simeq k$ canon.\unsure{t}).

\deleted{\unsure{O}}\unsure{outre} les \unsure{données} de commutativité, d'associativité (à vérifier leurs compatibilités !), a-t-on une unité ?

Il faudrait pour le moins, \deleted{que} pour une unité,
\diagram{\begin{tikzcd}[ampersand replacement=\&]0\arrow[r,dash]\&E^\circ\arrow[r,dash]\&E\arrow[r,dash]\&F\arrow[r]\&0\end{tikzcd}}
que $E^\circ$, et $F$ soient unités pour le $\otimes$ ordinaire, donc \deleted{inversibles} \inserted{$\simeq k$}.
\diagram{\begin{tikzcd}[ampersand replacement=\&]0\arrow[r]\&k\arrow[r]\&E\arrow[r]\&k\arrow[r]\&0\end{tikzcd}}
\deleted{\unsure{Donc les}} \unsure{Choisissons} scindage, alors la $\lambda$-triv. est définie par
\[
\deleted{\unsure{k\to k}}:\ k\longrightarrow k\qquad\text{i.e. }\beta_0\in k
\]
J'ai envie de prendre $\beta_0=1$. Prenons plus généralement $\beta$ quelconque, je dis qu'alors
\[
E_\beta\mathbin{*}_{\lambda}E_2\simeq {}^\beta E_2,
\]
muni de la $\lambda$ \inserted{trivialisation $\beta u_2$ déduite de celle de $E_2$}
\margin{179}
